\clearpage
\phantomsection
\addcontentsline{toc}{chapter}{Lexique}
\pdfbookmark[0]{Lexique}{lexique}
\origpage{465--468}
\markboth{Lexique}{Lexique}
\thispagestyle{plain}

\begin{center}
{\Huge\itshape\rmfamily\color{BellaLavenderDeep} Lexique}
\end{center}
\vspace{1.2em}

\begingroup
\small
\setlength{\parindent}{0pt}
\setlength{\parskip}{0.18em}
\raggedcolumns
\begin{multicols}{2}

abélien, 4.3

adjoint, 12.40

Adjonction symbolique, 7.62

\BellaCorrectionNote{d'Alembert--Gauss (théorème de), 7.71}{Alembert Gauss, (théorème d'), 7.71}{Le nom usuel du théorème est d'Alembert--Gauss.}

algèbre, 6.29

algèbre extérieure, 9.9

algébrique (élément), 7.57

algébrique (extension), 7.68

algébriquement clos, 7.51

alternée (forme), 9.2

alterné (groupe), 4.55

antisymétrique (loi), 2.2

antisymétrique (forme), 9.3

antisymétrisée, 9.6

appartenance, 1.3

application, 1.12

associativité, 1.24, 2.26

automorphisme, 2.22, 6.25

\medskip
base, 6.60

base incomplète (théorème de la), 6.66 et 6.72

Bézout (théorème de), 7.30

bijective, 1.17

bilinéaire, 11.1

binaire (relation), 2.1

bon (ordre), 2.43

borne (supérieure, inférieure), 2.20

borné, 2.20

\medskip
Cayley--Hamilton (théorème de), 10.44

caractéristique (d'un anneau), 5.29

caractéristique (polynôme), 10.19

caractéristiques (sous-espaces), 10.53

cardinaux (axiome des), 3.2

\BellaCorrectionNote{Cauchy--Schwarz (théorème de), 11.57, 12.32}{Cauchy Schwartz (théorème de), 11.57, 12.32}{Le nom usuel de l'inégalité est Cauchy--Schwarz.}

centre, 4.71

chaîne, 2.54

choix (axiome du), 2.45

classe d'équivalence, 2.5

clos (algébriquement), 7.51

clôture algébrique, 7.69

codimension, 6.91

coefficient directeur, 7.14

cofacteur, 9.32

comatrice, 9.40

combinaison linéaire, 6.18

commutative, 2.25

compatible, 2.21, 2.38

complémentaire (ensembliste), 1.7

composition (interne), 2.24

cône isotrope, 11.11

congruentes, 12.15

conjugués, 11.9, 12.18

conjugué (d'une partie), 11.14, 12.18

coordonnées, 6.76

corps, 5.13

Cramer (système de), 9.55

\medskip
cyclique, 4.13

\medskip
décomposition (corps de), 7.65

définie (forme), 11.12, 12.24

dégénérée (forme), 11.16, 12.22

degré, 7.2

demi-groupe, 4.26

dépendants, 6.57

déterminant, 9.14

diagonale (matrice), 8.13

diagonalisable, 10.9

discriminant, 11.42

distingué (sous-groupe), 4.18

distributive, 1.26, 2.33

diviseur de zéro, 5.8

domaine d'intégrité, 5.9

dual, 6.89

Dunford (théorème de), 10.54

\medskip
égalité, 1.1

endomorphisme, 2.22, 6.25

entier, entier naturel, 3.30

équipotent, 3.1

équivalence, 2.3

équivalence d'application, 2.39

équivalentes (matrices), 8.24

\BellaCorrectionNote{Erdős--Kaplansky (théorème d'), 6.113}{Erdös Kaplansky (théorème d'), 6.113}{Orthographe normalisée des noms Erdős et Kaplansky.}

espace vectoriel, 6.1

euclidien (anneau), 5.27

euclidienne (division), 7.19

\BellaCorrectionNote{Euler (identité d'), 11.47}{Euler (identité d'), 11.4 7}{Le numéro est 11.47.}

extension, 7.54

externe, 2.34

extrémal, 7.21

\medskip
factorisation des morphismes, 2.40

Fermat (théorème de), 7.13

fini (cardinal), 3.30

fini (ensemble), 3.34

fonction, (1.11)

forme linéaire, (6.89)

\medskip
Gauss (entiers de), exercice 8, chapitre 5

Gauss (méthode de), 11.53

Gauss (théorème de), 5.52

génératrice, 6.59

graphe, 1.11

groupe, 4.1

groupe linéaire, 6.30

\medskip
Hadamard (localisation des valeurs propres), 10.22

héréditaire (partie), 2.50

hermitien (espace), 12.36

hermitien (opérateur), 12.55

hermitienne (forme), 12.5

hermitienne (matrice), 12.13

Hilbert (espace de), 12.36

homothétie, 10.2

homomorphisme, 6.25

\medskip
idéal, 5.16

idempotent, 2.27

image (directe, réciproque), 1.13

inclusion, 1.5

indexation, 1.14

indice (d'un sous-groupe), 4.60

inductif, 2.41

infini (axiome de l'), 3.42

infini (cardinal), 3.30

injective, 1.15

intègre, 5.9

interne (loi de composition), 2.24

intersection, 1.23

irréductible, 7.21

isomorphisme, 2.22, 6.25

isomorphisme (premier théorème), 4.23, 6.41

isomorphisme (deuxième théorème), 6.43

isotrope, 11.10, 12.24

isotropie (groupe d'), 4.73

\medskip
\BellaCorrectionNote{Jordan (réduite élémentaire), 10.59}{Jordan (réduite élémentaire), 10,59}{Le séparateur du numéro est un point.}

Jordan (réduite de), 10.61

\medskip
majorant, 2.19

majoré, 2.19

matrice, 8.3

maximal (élément), 2.12

maximal (idéal), 5.42

maximum, 2.13

mineur, 9.34

minimal (élément), 2.17

minimal (polynôme), 10.43

minimum, 2.18

Minkowski (inégalité de), 11.61, 12.35

minorant, 2.19

minoré, 2.19

monogène, 4.13

morphisme, 2.21, 2.35

multiplicité (d'un zéro), 7.39

\medskip
négative (forme quadratique), 11.55

neutre, 2.28

nilpotent, 10.56

normal, 12.65

noyau (forme quadratique), 11.16

noyau (forme hermitienne), 12.22

noyau, 4.17

noyau (théorème des), 10.47

\medskip
opérant (groupe), 4.63

opérateur, 10.4

orbite, 4.65

ordre (d'un groupe), 4.58

ordre (d'un zéro), 7.39

ordre (bon), 2.43

ordre (relation), 2.7

orthogonal (groupe), 11.28

orthogonal (opérateur), 11.26

orthogonal (d'une partie), 6.94

orthogonalité (dual), 6.93

orthonormalisation de Schmidt, 12.37

\medskip
partition, 2.6

$p$-linéaire, 9.1

p.g.c.d., 7.26

plongement, 5.33

polaire (forme), 11.8, 12.7

polynôme, 6.13, 7.1

polynôme quadratique, 11.46

positive (f. quadratique), 11.55

p.p.c.m., 7.36

préhilbertien, 12.36

premier, 7.21

premier (idéal), 5.46

premier (nombre), 5.53

premiers entre eux, 7.28

presque tous nuls, 6.75bis

principal (anneau), 5.26

principal (idéal), 5.19

produit de composition, 1.18

produit extérieur, 9.7

\BellaCorrectionNote{projecteur, 6.47}{projecteur, 6.4 7}{Le numéro est 6.47.}

propre (valeur), 10.4

propre (vecteur), 10.6

propre (sous-espace), 10.6

\medskip
quadratique (forme), 11.4

quotient (ensemble), 2.5

\medskip
racine, 7.37

rang (application linéaire), 6.86

rang (famille de vecteurs), 9.44

rang (matrice), 8.26

rang (forme quadratique), 11.33

réciproque (application), 1.21

réflexive, 2.2

relation, 2.1

régulier, 2.36, 9.43

réunion, 1.22

Rouché (théorème de), 9.51

rupture (corps de), 7.64

\medskip
scindé, 10.28

segment (dans $E$ ordonné), 3.16

semblables, 8.25

semi-linéaire, 12.1

sesquilinéaire, 12.2

sesquilinéaire à symétrie hermitienne, 12.4

signature (forme quadratique), 11.64

signature (permutation), 4.52

singleton, 1.10

somme directe, 6.45

sous-espace vectoriel, 6.9

sous-groupe, 4.4

spectrale (valeur), 10.4

spectre, 10.5

stabilisateur, 4.73

Steinitz (théorème de), 7.70

strict (ordre), 2.11

suite, 6.11

supplémentaires, 6.46

surjective, 1.16

Sylvester (théorème de), 11.63

symétrique (relation), 2.2

symétrique (forme), 11.3

symétrique (groupe), 4.46

symétrisé, 4.28

\medskip
Taylor, 7.44

trace, 10.17

trajectoire, 4.65

transcendant, 7.56

transitive, 2.2

transport de structure, 2.23

transposée (matrice), 8.8

transposition, 4.48, 6.100

trigonalisable, 10.30

\medskip
unitaire (anneau), 5.4

unitaire (groupe), 12.48

unitaire (matrice), 12.50

unitaire (opérateur), 12.45

unitaire (polynôme), 7.15

unitaire spécial, 12.54

\medskip
valuation (facteur premier), 7.34

valuation (polynôme), 7.2

\BellaCorrectionNote{Vandermonde (exercice 6 chapitre 9)}{Van der Mond (exercice 6 chapitre 9)}{Le déterminant classique porte le nom de Vandermonde.}

vide, 1.9

\medskip
zéro, 7.37

Zermelo (axiome de), 2.44

Zorn (axiome de), 2.42

\end{multicols}

\vfill
\begin{center}
\footnotesize\sffamily\color{BellaGray}
Imprimé en France\\
Imprimerie des Presses Universitaires de France\\
73, avenue Ronsard, 41100 Vendôme\\
Août 1993 -- N$^{\circ}$ 39 719
\end{center}
\endgroup
